\section{Expected Outcome}

Based on the stated objectives, FeatureSAGE is expected to improve the SAGE pipeline by replacing the pSp-centered inversion stage with the StyleFeatureEditor inverter and Feature Editor. The expected improvement is not a change in the underlying StyleGAN2 generator or in the attribute group editing objective; rather, it is an improvement in how input images are represented before editing. By extracting both $W^+$ latent codes and intermediate feature tensors, the proposed framework is expected to preserve more local detail while retaining the ability to apply category-irrelevant edits.

The expected outcomes are evaluated on the FUNIT Animal Faces and 102 Flowers datasets under the same one-shot protocol used for the baseline methods. FID is used to measure how closely generated images align with real-image distributions, while LPIPS is used to examine perceptual variation among generated outputs. Within this evaluation design, the expected measurable outcomes are:

\begin{enumerate}
    \item \textbf{Lower FID}, indicating improved visual realism and stronger distributional alignment between real and synthetic images.

    \item \textbf{Comparable or higher LPIPS as a diversity measure}, indicating that improved fidelity does not eliminate perceptual variation among generated samples.
\end{enumerate}

Together, these outcomes test whether the FeatureSAGE modifications improve few-shot image generation while preserving the edit-based structure of the original SAGE framework.
